\chapter{Despliegue}

\section{Requisitos del Servidor}

\begin{itemize}
    \item Java 21 o superior (OpenJDK 21 LTS recomendado).
    \item PostgreSQL 14 o superior.
    \item Maven 3.8+ (para compilaci\'on desde c\'odigo fuente).
    \item Al menos 512 MB de RAM para la aplicaci\'on (1 GB recomendado).
    \item Conexi\'on SMTP saliente para notificaciones por email (opcional).
    \item Cuenta de Stripe con API keys para pagos en l\'inea (opcional).
\end{itemize}

\section{Configuraci\'on de Base de Datos}

\subsection{Crear Base de Datos PostgreSQL}

\begin{lstlisting}[caption=Creaci\'on de base de datos y usuario en PostgreSQL]
sudo -u postgres psql -c "CREATE DATABASE veterinary_db;"
sudo -u postgres psql -c "CREATE USER veterinary_user WITH PASSWORD 'veterinary_pass';"
sudo -u postgres psql -c "GRANT ALL PRIVILEGES ON DATABASE veterinary_db TO veterinary_user;"
\end{lstlisting}

\subsection{Verificar Conexi\'on}

\begin{lstlisting}[caption=Verificar que la base de datos es accesible]
psql -U veterinary_user -d veterinary_db -h localhost
\end{lstlisting}

\section{Ejecuci\'on en Desarrollo}

\begin{lstlisting}[caption=Ejecutar la aplicaci\'on en modo desarrollo]
mvn spring-boot:run
\end{lstlisting}

La aplicaci\'on estar\'a disponible en \texttt{http://localhost:8090}. La base de datos se crea y puebla autom\'aticamente al primer inicio gracias a \texttt{spring.jpa.hibernate.ddl-auto=update} y el \texttt{DataInitializer}.

\section{Compilaci\'on para Producci\'on}

\begin{lstlisting}[caption=Compilar y empaquetar como JAR]
mvn clean package -DskipTests
\end{lstlisting}

Genera el archivo \texttt{target/veterinary-clinic-management-system-1.0.0.jar}.

\begin{lstlisting}[caption=Ejecutar JAR compilado con variables de entorno]
export JWT_SECRET="clave-secreta-muy-segura-12345"
export JWT_EXPIRATION=1800000
export MAIL_HOST=smtp.gmail.com
export MAIL_USERNAME=notificaciones@petclinic.com
export MAIL_PASSWORD=contrasena-de-app
export STRIPE_API_KEY=sk_live_...
export STRIPE_WEBHOOK_SECRET=whsec_...

java -jar target/veterinary-clinic-management-system-1.0.0.jar
\end{lstlisting}

\section{Despliegue con Docker}

\subsection{Dockerfile}

\begin{lstlisting}[caption=Dockerfile multi-etapa para imagen optimizada]
FROM maven:3.9-eclipse-temurin-21 AS builder
WORKDIR /app
COPY pom.xml .
RUN mvn dependency:go-offline
COPY src ./src
RUN mvn clean package -DskipTests

FROM openjdk:21-jdk-slim
WORKDIR /app
COPY --from=builder /app/target/*.jar app.jar
EXPOSE 8090
ENTRYPOINT ["java", "-jar", "/app/app.jar"]
\end{lstlisting}

\subsection{Docker Compose}

\begin{lstlisting}[caption=Docker Compose con PostgreSQL y aplicaci\'on]
version: '3.8'
services:
  postgres:
    image: postgres:14
    environment:
      POSTGRES_DB: veterinary_db
      POSTGRES_USER: veterinary_user
      POSTGRES_PASSWORD: veterinary_pass
    ports:
      - "5432:5432"
    volumes:
      - postgres_data:/var/lib/postgresql/data

  app:
    build: .
    ports:
      - "8090:8090"
    environment:
      SPRING_DATASOURCE_URL: jdbc:postgresql://postgres:5432/veterinary_db
      SPRING_DATASOURCE_USERNAME: veterinary_user
      SPRING_DATASOURCE_PASSWORD: veterinary_pass
      JWT_SECRET: clave-secreta-docker-2024
      JWT_EXPIRATION: 1800000
    depends_on:
      - postgres

volumes:
  postgres_data:
\end{lstlisting}

\begin{lstlisting}[caption=Construir y ejecutar con Docker Compose]
docker-compose build
docker-compose up -d
docker-compose logs -f app  # Ver logs en tiempo real
\end{lstlisting}

\section{Despliegue con systemd (Producci\'on Linux)}

Para ejecutar la aplicaci\'on como servicio del sistema:

\begin{lstlisting}[caption=Archivo de servicio systemd]
[Unit]
Description=PetClinic Veterinary Clinic
After=network.target postgresql.service

[Service]
Type=simple
User=petclinic
WorkingDirectory=/opt/petclinic
Environment="JWT_SECRET=..."
Environment="JWT_EXPIRATION=1800000"
Environment="MAIL_HOST=smtp.gmail.com"
ExecStart=/usr/bin/java -jar /opt/petclinic/app.jar
Restart=on-failure
RestartSec=10

[Install]
WantedBy=multi-user.target
\end{lstlisting}

\begin{lstlisting}[caption=Instalar y gestionar el servicio]
sudo cp petclinic.service /etc/systemd/system/
sudo systemctl daemon-reload
sudo systemctl enable petclinic
sudo systemctl start petclinic
sudo systemctl status petclinic
\end{lstlisting}

\section{Proxy Inverso con Nginx}

Para entornos de producci\'on con dominio propio:

\begin{lstlisting}[caption=Configuraci\'on de Nginx como proxy inverso]
server {
    listen 80;
    server_name petclinic.midominio.com;

    location / {
        proxy_pass http://localhost:8090;
        proxy_set_header Host $host;
        proxy_set_header X-Real-IP $remote_addr;
        proxy_set_header X-Forwarded-For $proxy_add_x_forwarded_for;
        proxy_set_header X-Forwarded-Proto $scheme;
    }

    location /ws {
        proxy_pass http://localhost:8090;
        proxy_http_version 1.1;
        proxy_set_header Upgrade $http_upgrade;
        proxy_set_header Connection "upgrade";
        proxy_set_header Host $host;
    }
}
\end{lstlisting}

\section{Variables de Entorno}

El sistema soporta configuraci\'on mediante variables de entorno para entornos productivos, lo que permite despliegues sin modificar archivos de configuraci\'on:

\begin{longtable}{|p{4cm}|p{6cm}|p{4cm}|}
\hline
\textbf{Variable} & \textbf{Descripci\'on} & \textbf{Valor por Defecto} \\
\hline
JWT\_SECRET & Clave secreta para firmar tokens JWT (HMAC-SHA384) & (valor interno) \\
JWT\_EXPIRATION & Duraci\'on del token en milisegundos & 1800000 (30 min) \\
MAIL\_HOST & Servidor SMTP para env\'io de correos & smtp.gmail.com \\
MAIL\_PORT & Puerto SMTP & 587 \\
MAIL\_USERNAME & Usuario SMTP & (vac\'io) \\
MAIL\_PASSWORD & Contrase\~na SMTP & (vac\'io) \\
STRIPE\_API\_KEY & Clave API de Stripe (pk\_/sk\_) & sk\_test\_placeholder \\
STRIPE\_WEBHOOK\_SECRET & Secreto de webhook Stripe & whsec\_placeholder \\
FRONTEND\_URL & URL del frontend para enlaces en emails & http://localhost:8090 \\
UPLOAD\_DIR & Directorio de subida de archivos & ./uploads \\
\hline
\end{longtable}

\section{Estructura Completa del Proyecto}

\begin{verbatim}
veterinary-clinic-management-system/
  |-- pom.xml                          # Dependencias Maven
  |-- README.md                        # Instrucciones rapidas
  |-- .gitignore
  |-- Dockerfile                       # Imagen Docker multi-etapa
  |-- docker-compose.yml               # Orquestacion Docker
  |-- docs/                            # Documentacion LaTeX
  |   |-- main.tex
  |   |-- chapters/
  |   |-- images/
  |-- src/
      |-- main/
      |   |-- java/com/veterinary/
      |   |   |-- VeterinaryApplication.java
      |   |   |-- config/
      |   |   |-- controller/          # 20 REST controllers
      |   |   |-- domain/              # 20 entidades + 9 enums
      |   |   |-- dto/                 # 6 DTOs
      |   |   |-- repository/          # 19 repositorios
      |   |   |-- scheduler/
      |   |   |-- security/            # 4 clases de seguridad
      |   |   |-- service/             # 22 servicios
      |   |-- resources/
      |       |-- application.properties
      |       |-- static/
      |           |-- index.html
      |           |-- css/app.css
      |           |-- js/{api.js, app.js, vet-extras.js}
      |-- test/
          |-- java/com/veterinary/
          |   |-- VeterinarySystemIntegrationTest.java   # 15 tests
          |   |-- NewFeaturesIntegrationTest.java         # 15 tests
          |-- resources/
              |-- application.properties
\end{verbatim}

\section{Comandos \'Utiles}

\begin{longtable}{|p{6cm}|p{8cm}|}
\hline
\textbf{Comando} & \textbf{Descripci\'on} \\
\hline
mvn spring-boot:run & Ejecutar la aplicaci\'on en desarrollo \\
mvn clean install & Compilar y empaquetar \\
mvn test & Ejecutar pruebas unitarias y de integraci\'on \\
mvn test -Dtest=ClaseTest\#metodo & Ejecutar prueba espec\'ifica \\
docker-compose up -d & Desplegar con Docker Compose \\
docker-compose logs -f app & Ver logs de la aplicaci\'on \\
java -jar target/*.jar & Ejecutar JAR compilado \\
sudo systemctl start petclinic & Iniciar servicio systemd \\
\hline
\end{longtable}

\section{Resoluci\'on de Problemas}

\subsection{Error de Conexi\'on a Base de Datos}

Verificar que PostgreSQL est\'e en ejecuci\'on y accesible:

\begin{lstlisting}[caption=Verificar conexi\'on PostgreSQL]
sudo systemctl status postgresql
psql -U veterinary_user -d veterinary_db -h localhost
\end{lstlisting}

Si el error persiste, verificar que \texttt{pg\_hba.conf} permita conexiones locales con autenticaci\'on por contrase\~na.

\subsection{Puerto Ocupado}

Si el puerto 8090 ya est\'a en uso, se puede cambiar en \texttt{application.properties} o v\'ia variable de entorno:

\begin{lstlisting}
server.port=8091
# O via entorno:
export SERVER_PORT=8091
\end{lstlisting}

\subsection{Error de Compilaci\'on}

Asegurar que Java 21 est\'e configurado correctamente:

\begin{lstlisting}
java -version  # Debe mostrar 21.x
mvn -version   # Debe usar Java 21
\end{lstlisting}

\subsection{Error de JWT}

Si los tokens JWT no se validan correctamente, verificar que \texttt{JWT\_SECRET} tenga al menos 256 bits (32 caracteres) y sea consistente entre reinicios de la aplicaci\'on.

\subsection{WebSocket no conecta}

Verificar que el proxy inverso (si existe) tenga configurada la actualizaci\'on de protocolo WebSocket (\texttt{Upgrade} y \texttt{Connection} headers). Para Nginx, asegurar que la ubicaci\'on \texttt{/ws} tenga las directivas \texttt{proxy\_set\_header Upgrade} y \texttt{proxy\_set\_header Connection "upgrade"}.
